\renewcommand{\Anno}{2007}

\renewcommand{\NumeroProblema}{1}

\renewcommand{\Testo}{Sono dati un quadrato di lato $1$ ed un
quadrato di lato $\ell$.  Il centro del quadrato di lato $\ell$ coincide con uno dei vertici del quadrato di lato 1.

\begin{center}
	\psset{unit=7ex}
	\pspicture(-1,-1)(2.5,2.3)
	\SpecialCoor
	\psclip{\pscustom{\translate(1,1)\rotate{60}\psframe(-1,-1)(1,1)}}
	\psframe[fillstyle=solid,fillcolor=cyan,linestyle=none](-1,-1)(1,1)
	\endpsclip
	\psframe(-1,-1)(1,1)
	\pscustom{\translate(1,1)\rotate{60}\psframe(-1,-1)(1,1)}
	\endpspicture
\end{center}

	\begin{enumerate}
	    \item  Dimostrare che, nel caso particolare $\ell=1$, l'area dell'intersezione tra i 2 quadrati non dipende dall'angolo di rotazione.
	
	    \item  Determinare se la stessa proprietà resta vera per tutti i valori $\ell>0$.
	\end{enumerate}
}

\renewcommand{\Risposta}{
\begin{enumerate}
	\item  \ldots
		
	\item  Non resta vera
\end{enumerate}
}

\renewcommand{\Aiutino}{
\begin{enumerate}
	\renewcommand{\labelenumi}{(\alph{enumi})} 
	\item Considerare la seguente figura: 
	\begin{center}
	\psset{unit=6ex}
	\pspicture(-1,-1)(2.5,2.5)
	\def\angolo{60}
	\def\lato{1}
	\SpecialCoor
	\pnode(-1,-1){A}
	\pnode(-1,1){B}
	\pstGeonode[PointName=A,PosAngle=45,PointSymbol=none](1,1){C}
	\pnode(1,-1){D}
	\pnode(0,0){O}
	\pstRotation[RotAngle=\angolo,PointSymbol=none,PointName=none]{O}{A}[AA]
	\pstRotation[RotAngle=\angolo,PointSymbol=none,PointName=none]{O}{B}[BB]
	\pstRotation[RotAngle=\angolo,PointSymbol=none,PointName=none]{O}{C}[CC]
	\pstRotation[RotAngle=\angolo,PointSymbol=none,PointName=none]{O}{D}[DD]
	\pstHomO[HomCoef=\lato,PointSymbol=none,PointName=none]{O}{AA}[AAA]
	\pstHomO[HomCoef=\lato,PointSymbol=none,PointName=none]{O}{BB}[BBB]
	\pstHomO[HomCoef=\lato,PointSymbol=none,PointName=none]{O}{CC}[CCC]
	\pstHomO[HomCoef=\lato,PointSymbol=none,PointName=none]{O}{DD}[DDD]
	\pstTranslation[PointName=B',PosAngle=-90,PointSymbol=none,PointName=none]{O}{C}{AAA}[A']
	\pstTranslation[PointName=A',PosAngle=180,PointSymbol=none,PointName=none]{O}{C}{BBB}[B']
	\pstTranslation[PointSymbol=none,PointName=none]{O}{C}{CCC}[C']
	\pstTranslation[PointSymbol=none,PointName=none]{O}{C}{DDD}[D']
	\pstInterLL[PosAngle=-135,PointSymbol=none,PointName=none]{A'}{B'}{D}{C}{K}
	\pstInterLL[PosAngle=135,PointSymbol=none,PointName=none]{C'}{B'}{B}{C}{K'}
	\psclip{\psline(A')(B')(C')(D')(A')}
	\psframe[fillstyle=solid,fillcolor=cyan,linestyle=none](A)(C)
	\endpsclip
	\psline[fillstyle=solid,fillcolor=green,linestyle=none](C)(B')(K')
	\psline(C)(B')
	\psline(C)(A')
	\psframe(A)(C)
	\psline(A')(B')(C')(D')(A')
% 	\pstMarkAngle{K}{C}{A'}{}
% 	\pstMarkAngle{K'}{C}{B'}{}
	\endpspicture
	\hspace{6em}
	\pspicture(-1,-1)(2.5,2.5)
	\def\angolo{60}
	\def\lato{1}
	\SpecialCoor
	\pnode(-1,-1){A}
	\pnode(-1,1){B}
	\pstGeonode[PointName=A,PosAngle=45,PointSymbol=none](1,1){C}
	\pnode(1,-1){D}
	\pnode(0,0){O}
	\pstRotation[RotAngle=\angolo,PointSymbol=none]{O}{A}[AA]
	\pstRotation[RotAngle=\angolo,PointSymbol=none]{O}{B}[BB]
	\pstRotation[RotAngle=\angolo,PointSymbol=none]{O}{C}[CC]
	\pstRotation[RotAngle=\angolo,PointSymbol=none]{O}{D}[DD]
	\pstHomO[HomCoef=\lato,PointSymbol=none]{O}{AA}[AAA]
	\pstHomO[HomCoef=\lato,PointSymbol=none]{O}{BB}[BBB]
	\pstHomO[HomCoef=\lato,PointSymbol=none]{O}{CC}[CCC]
	\pstHomO[HomCoef=\lato,PointSymbol=none]{O}{DD}[DDD]
	\pstTranslation[PointName=B',PosAngle=-90,PointSymbol=none]{O}{C}{AAA}[A']
	\pstTranslation[PointName=A',PosAngle=180,PointSymbol=none]{O}{C}{BBB}[B']
	\pstTranslation[PointSymbol=none]{O}{C}{CCC}[C']
	\pstTranslation[PointSymbol=none]{O}{C}{DDD}[D']
	\pstInterLL[PosAngle=-135,PointSymbol=none]{A'}{B'}{D}{C}{K}
	\pstInterLL[PosAngle=135,PointSymbol=none]{C'}{B'}{B}{C}{K'}
	\psclip{\psline(A')(B')(C')(D')(A')}
	\psframe[fillstyle=solid,fillcolor=cyan,linestyle=none](A)(C)
	\endpsclip
	\psline[fillstyle=solid,fillcolor=green,linestyle=none](C)(A')(K)
	\psline[fillstyle=solid,fillcolor=white,linestyle=none](C)(B')(K')
	\psline(C)(B')
	\psline(C)(A')
	\psframe(A)(C)
	\psline(A')(B')(C')(D')(A')
% 	\pstMarkAngle{K}{C}{A'}{}
% 	\pstMarkAngle{K'}{C}{B'}{}
	\endpspicture
\end{center}


	\item  Considerare il caso $\ell=2$.
\end{enumerate}
}

\renewcommand{\SoluzioneUfficiale}{\TitoloSoluzione \DomandaA
Mostriamo che l'area dell'intersezione è uguale ad $1/4$ dell'area del quadrato di lato $\ell$.  Con riferimento alla figura sottostante ci basta dimostrare che il triangolo $AA'K'$ è uguale al triangolo $AB'K$.
\begin{center}
	\psset{unit=8ex}
	\pspicture(-1,-1)(2.5,2.5)
	\def\angolo{60}
	\def\lato{1}
	\SpecialCoor
	\pnode(-1,-1){A}
	\pnode(-1,1){B}
	\pstGeonode[PointName=A,PosAngle=45](1,1){C}
	\pnode(1,-1){D}
	\pnode(0,0){O}
	\pstRotation[RotAngle=\angolo,PointSymbol=none]{O}{A}[AA]
	\pstRotation[RotAngle=\angolo,PointSymbol=none]{O}{B}[BB]
	\pstRotation[RotAngle=\angolo,PointSymbol=none]{O}{C}[CC]
	\pstRotation[RotAngle=\angolo,PointSymbol=none]{O}{D}[DD]
	\pstHomO[HomCoef=\lato,PointSymbol=none]{O}{AA}[AAA]
	\pstHomO[HomCoef=\lato,PointSymbol=none]{O}{BB}[BBB]
	\pstHomO[HomCoef=\lato,PointSymbol=none]{O}{CC}[CCC]
	\pstHomO[HomCoef=\lato,PointSymbol=none]{O}{DD}[DDD]
	\pstTranslation[PointName=B',PosAngle=-90]{O}{C}{AAA}[A']
	\pstTranslation[PointName=A',PosAngle=180]{O}{C}{BBB}[B']
	\pstTranslation[PointSymbol=none]{O}{C}{CCC}[C']
	\pstTranslation[PointSymbol=none]{O}{C}{DDD}[D']
	%\pstProjection{A'}{B'}{C}{H}
	%\pstProjection{B'}{C'}{C}{H'}
	\pstInterLL[PosAngle=-135]{A'}{B'}{D}{C}{K}
	\pstInterLL[PosAngle=135]{C'}{B'}{B}{C}{K'}
	\psclip{\psline(A')(B')(C')(D')(A')}
	\psframe[fillstyle=solid,fillcolor=cyan,linestyle=none](A)(C)
	\endpsclip
	\psline[fillstyle=solid,fillcolor=green,linestyle=none](C)(B')(K')
	\psline(C)(B')
	\psline(C)(A')
	\psframe(A)(C)
	\psline(A')(B')(C')(D')(A')
	\pstMarkAngle{K}{C}{A'}{}
	\pstMarkAngle{K'}{C}{B'}{}
	\endpspicture
	\hspace{6em}
	\pspicture(-1,-1)(2.5,2.5)
	\def\angolo{60}
	\def\lato{1}
	\SpecialCoor
	\pnode(-1,-1){A}
	\pnode(-1,1){B}
	\pstGeonode[PointName=A,PosAngle=45](1,1){C}
	\pnode(1,-1){D}
	\pnode(0,0){O}
	\pstRotation[RotAngle=\angolo,PointSymbol=none]{O}{A}[AA]
	\pstRotation[RotAngle=\angolo,PointSymbol=none]{O}{B}[BB]
	\pstRotation[RotAngle=\angolo,PointSymbol=none]{O}{C}[CC]
	\pstRotation[RotAngle=\angolo,PointSymbol=none]{O}{D}[DD]
	\pstHomO[HomCoef=\lato,PointSymbol=none]{O}{AA}[AAA]
	\pstHomO[HomCoef=\lato,PointSymbol=none]{O}{BB}[BBB]
	\pstHomO[HomCoef=\lato,PointSymbol=none]{O}{CC}[CCC]
	\pstHomO[HomCoef=\lato,PointSymbol=none]{O}{DD}[DDD]
	\pstTranslation[PointName=B',PosAngle=-90]{O}{C}{AAA}[A']
	\pstTranslation[PointName=A',PosAngle=180]{O}{C}{BBB}[B']
	\pstTranslation[PointSymbol=none]{O}{C}{CCC}[C']
	\pstTranslation[PointSymbol=none]{O}{C}{DDD}[D']
	%\pstProjection{A'}{B'}{C}{H}
	%\pstProjection{B'}{C'}{C}{H'}
	\pstInterLL[PosAngle=-135]{A'}{B'}{D}{C}{K}
	\pstInterLL[PosAngle=135]{C'}{B'}{B}{C}{K'}
	\psclip{\psline(A')(B')(C')(D')(A')}
	\psframe[fillstyle=solid,fillcolor=cyan,linestyle=none](A)(C)
	\endpsclip
	\psline[fillstyle=solid,fillcolor=green,linestyle=none](C)(A')(K)
	\psline[fillstyle=solid,fillcolor=white,linestyle=none](C)(B')(K')
	\psline(C)(B')
	\psline(C)(A')
	\psframe(A)(C)
	\psline(A')(B')(C')(D')(A')
	\pstMarkAngle{K}{C}{A'}{}
	\pstMarkAngle{K'}{C}{B'}{}
	\endpspicture
\end{center}

Tale uguaglianza di triangoli segue dal secondo criterio di congruenza, in quanto $AA'=AB'$ (semidiagonali del quadrato di lato $\ell$), gli angoli $\angle AA'K'$ e $\angle AB'K$ sono uguali
(entrambi misurano $45^{\circ}$), e gli angoli $\angle A'AK'$ e $\angle
B'AK$ sono uguali (essendo entrambi complementari di $\angle A'AK$).

In alternativa si può osservare che la rotazione di $90^{\circ}$ in senso antiorario con centro in $A$ manda il triangolo $AA'K'$ nel triangolo $AB'K$.


\DomandaB Lo stesso ragionamento del punto (a) continua a valere
finché un vertice del quadrato di lato $\ell$, diciamo $A'$, cade all'interno del quadrato di lato 1.  Quando $\ell>\sqrt{2}$ ci sono invece posizioni per cui questo non accade, ed è facile convincersi che l'area dell'intersezione dipende dall'angolazione.

Ad esempio nella situazione mostrata in figura, in cui $\ell=2$, si ha che quando i lati dei 2 quadrati sono paralleli l'area dell'intersezione è pari a tutto il quadrato piccolo, mentre ruotando di $45^{\circ}$ l'area dell'intersezione è strettamente minore di quella del quadrato piccolo.

\begin{center}
	\psset{unit=5ex}
	\pspicture(-1,-1.3)(2.5,3.5)
	\def\angolo{0}
	\def\lato{2}
	\SpecialCoor
	\pnode(-1,-1){A}
	\pnode(-1,1){B}
	\pnode(1,1){C}
	\pnode(1,-1){D}
	\pnode(0,0){O}
	\pstRotation[RotAngle=\angolo,PointSymbol=none]{O}{A}[AA]
	\pstRotation[RotAngle=\angolo,PointSymbol=none]{O}{B}[BB]
	\pstRotation[RotAngle=\angolo,PointSymbol=none]{O}{C}[CC]
	\pstRotation[RotAngle=\angolo,PointSymbol=none]{O}{D}[DD]
	\pstHomO[HomCoef=\lato,PointSymbol=none]{O}{AA}[AAA]
	\pstHomO[HomCoef=\lato,PointSymbol=none]{O}{BB}[BBB]
	\pstHomO[HomCoef=\lato,PointSymbol=none]{O}{CC}[CCC]
	\pstHomO[HomCoef=\lato,PointSymbol=none]{O}{DD}[DDD]
	\pstTranslation[PointSymbol=none]{O}{C}{AAA}[A']
	\pstTranslation[PointSymbol=none]{O}{C}{BBB}[B']
	\pstTranslation[PointSymbol=none]{O}{C}{CCC}[C']
	\pstTranslation[PointSymbol=none]{O}{C}{DDD}[D']
	\psclip{\psline(A')(B')(C')(D')(A')}
	\psframe[fillstyle=solid,fillcolor=cyan,linestyle=none](A)(C)
	\endpsclip
	\psframe(A)(C)
	\psline(A')(B')(C')(D')(A')
	\endpspicture
	\hspace{10em}
	\pspicture(-1,-1.3)(2.5,3.5)
	\def\angolo{45}
	\def\lato{1.8}
	\SpecialCoor
	\pnode(-1,-1){A}
	\pnode(-1,1){B}
	\pnode(1,1){C}
	\pnode(1,-1){D}
	\pnode(0,0){O}
	\pstRotation[RotAngle=\angolo,PointSymbol=none]{O}{A}[AA]
	\pstRotation[RotAngle=\angolo,PointSymbol=none]{O}{B}[BB]
	\pstRotation[RotAngle=\angolo,PointSymbol=none]{O}{C}[CC]
	\pstRotation[RotAngle=\angolo,PointSymbol=none]{O}{D}[DD]
	\pstHomO[HomCoef=\lato,PointSymbol=none]{O}{AA}[AAA]
	\pstHomO[HomCoef=\lato,PointSymbol=none]{O}{BB}[BBB]
	\pstHomO[HomCoef=\lato,PointSymbol=none]{O}{CC}[CCC]
	\pstHomO[HomCoef=\lato,PointSymbol=none]{O}{DD}[DDD]
	\pstTranslation[PointSymbol=none]{O}{C}{AAA}[A']
	\pstTranslation[PointSymbol=none]{O}{C}{BBB}[B']
	\pstTranslation[PointSymbol=none]{O}{C}{CCC}[C']
	\pstTranslation[PointSymbol=none]{O}{C}{DDD}[D']
% 	\psline[linestyle=none,fillstyle=solid,fillcolor=green](C)(A')(B')(C)
	\psclip{\psline(A')(B')(C')(D')(A')}
	\psframe[fillstyle=solid,fillcolor=cyan,linestyle=none](A)(C)
	\endpsclip
	\psframe(A)(C)
	\psline(A')(B')(C')(D')(A')
% 	\psline[linestyle=dashed](B)(B')
% 	\psline[linestyle=dashed](D)(A')
	\endpspicture
\end{center}

Lo stesso fenomeno di ha per tutti i valori di $\ell$ fino a $2\sqrt{2}$. Per $\ell\geq 2\sqrt{2}$ si ha invece che l'intersezione è sempre tutto il quadrato piccolo, dunque la sua area non dipende dall'angolo di rotazione.

In conclusione l'area dell'intersezione dipende dall'angolo di rotazione se e solo se $\sqrt{2}<\ell<2\sqrt{2}$.
}

\renewcommand{\SoluzioniAlternative}{\TitoloSoluzioneAlternativaA
Un altro modo di mostrare che l'area dell'intersezione è 1/4 dell'area del quadrato grande è basato sull'uguaglianza tra i triangoli $AHK$ e $AH'K'$ nella seguente figura.
\begin{center}
	\psset{unit=8ex}
	\pspicture(-1,-1)(2.5,2.5)
	\def\angolo{60}
	\def\lato{1}
	\SpecialCoor
	\pnode(-1,-1){A}
	\pnode(-1,1){B}
	\pstGeonode[PointName=A,PosAngle=45](1,1){C}
	\pnode(1,-1){D}
	\pnode(0,0){O}
	\pstRotation[RotAngle=\angolo,PointSymbol=none]{O}{A}[AA]
	\pstRotation[RotAngle=\angolo,PointSymbol=none]{O}{B}[BB]
	\pstRotation[RotAngle=\angolo,PointSymbol=none]{O}{C}[CC]
	\pstRotation[RotAngle=\angolo,PointSymbol=none]{O}{D}[DD]
	\pstHomO[HomCoef=\lato,PointSymbol=none]{O}{AA}[AAA]
	\pstHomO[HomCoef=\lato,PointSymbol=none]{O}{BB}[BBB]
	\pstHomO[HomCoef=\lato,PointSymbol=none]{O}{CC}[CCC]
	\pstHomO[HomCoef=\lato,PointSymbol=none]{O}{DD}[DDD]
	\pstTranslation[PointSymbol=none]{O}{C}{AAA}[A']
	\pstTranslation[PointName=A',PosAngle=180]{O}{C}{BBB}[B']
	\pstTranslation[PointSymbol=none]{O}{C}{CCC}[C']
	\pstTranslation[PointSymbol=none]{O}{C}{DDD}[D']
	\pstProjection{A'}{B'}{C}[H]
	\pstProjection{B'}{C'}{C}[H']
	\pstInterLL[PosAngle=-135]{A'}{B'}{D}{C}{K}
	\pstInterLL[PosAngle=135]{C'}{B'}{B}{C}{K'}
	\psclip{\psline(A')(B')(C')(D')(A')}
	\psframe[fillstyle=solid,fillcolor=cyan,linestyle=none](A)(C)
	\endpsclip
	\psline[fillstyle=solid,fillcolor=green,linestyle=none](C)(H)(K)
	\psframe(A)(C)
	\psline(A')(B')(C')(D')(A')
	\psline[linestyle=dashed](C)(H)
	\psline[linestyle=dashed](C)(H')
	\pstRightAngle{K}{H}{C}
	\pstRightAngle{K'}{H'}{C}
	\pstMarkAngle{H}{C}{K}{}
	\pstMarkAngle{H'}{C}{K'}{}
	\endpspicture
	\hspace{6em}
	\pspicture(-1,-1)(2.5,2.5)
	\def\angolo{60}
	\def\lato{1}
	\SpecialCoor
	\pnode(-1,-1){A}
	\pnode(-1,1){B}
	\pstGeonode[PointName=A,PosAngle=45](1,1){C}
	\pnode(1,-1){D}
	\pnode(0,0){O}
	\pstRotation[RotAngle=\angolo,PointSymbol=none]{O}{A}[AA]
	\pstRotation[RotAngle=\angolo,PointSymbol=none]{O}{B}[BB]
	\pstRotation[RotAngle=\angolo,PointSymbol=none]{O}{C}[CC]
	\pstRotation[RotAngle=\angolo,PointSymbol=none]{O}{D}[DD]
	\pstHomO[HomCoef=\lato,PointSymbol=none]{O}{AA}[AAA]
	\pstHomO[HomCoef=\lato,PointSymbol=none]{O}{BB}[BBB]
	\pstHomO[HomCoef=\lato,PointSymbol=none]{O}{CC}[CCC]
	\pstHomO[HomCoef=\lato,PointSymbol=none]{O}{DD}[DDD]
	\pstTranslation[PointSymbol=none]{O}{C}{AAA}[A']
	\pstTranslation[PointName=A',PosAngle=180]{O}{C}{BBB}[B']
	\pstTranslation[PointSymbol=none]{O}{C}{CCC}[C']
	\pstTranslation[PointSymbol=none]{O}{C}{DDD}[D']
	\pstProjection{A'}{B'}{C}[H]
	\pstProjection{B'}{C'}{C}[H']
	\pstInterLL[PosAngle=-135]{A'}{B'}{D}{C}{K}
	\pstInterLL[PosAngle=135]{C'}{B'}{B}{C}{K'}
	\psclip{\psline(A')(B')(C')(D')(A')}
	\psframe[fillstyle=solid,fillcolor=cyan,linestyle=none](A)(C)
	\endpsclip
	\psline[fillstyle=solid,fillcolor=green,linestyle=none](C)(H')(K')
	\psline[fillstyle=solid,fillcolor=white,linestyle=none](C)(H)(K)
	\psframe(A)(C)
	\psline(A')(B')(C')(D')(A')
	\psline[linestyle=dashed](C)(H)
	\psline[linestyle=dashed](C)(H')
	\pstRightAngle{K}{H}{C}
	\pstRightAngle{K'}{H'}{C}
	\pstMarkAngle{H}{C}{K}{}
	\pstMarkAngle{H'}{C}{K'}{}
	\endpspicture
\end{center}

Tale uguaglianza di triangoli segue dal fatto che sono entrambi triangoli rettangoli, che i cateti $AH$ e $AH'$ sono entrambi lunghi $\ell/2$, e che $\angle HAK=\angle H'AK'$ in quanto complementari dello stesso angolo $\angle HAK'$.

Come nel caso precedente, l'uguaglianza tra i triangoli $AHK$ e $AH'K'$ si può dimostrare anche con un semplice argomento di rotazione.

\TitoloSoluzioneAlternativaB Con riferimento alla figura iniziale,
indichiamo con $\theta$ l'angolo $\angle A'AK'$.  Nel triangolo $AA'K'$ conosciamo il lato $AA'$ ed i 2 angoli adiacenti.  Per una nota formula trigonometrica la sua area è
$$\mbox{Area}(AA'K')=\frac{1}{2}(AA')^{2}\cdot
\frac{\sin\theta\cdot\sin(45^{\circ})}{\sin(135^{\circ}-\theta)}.$$

Similmente nel triangolo $AA'K$ conosciamo il lato $AA'$ ed i 2 angoli adiacenti, dunque
$$\mbox{Area}(AA'K)=\frac{1}{2}(AA')^{2}\cdot
\frac{\sin(90^{\circ}-\theta)\cdot
\sin(45^{\circ})}{\sin(45^{\circ}+\theta)}.$$

Sommando queste 2 relazioni, e sfruttando che $\sin(135^{\circ}-\theta)=\sin(45^{\circ}+\theta)$ (in quanto angoli supplemetari) e la formula di addizione per calcolare $\sin(45^{\circ}+\theta)$, si vede che la dipendenza da $\theta$ scompare.
}

\renewcommand{\Osservazioni}{
% \TitoloOsservazione Si noti che il
% controesempio presentato nel punto (b) richiede che sia
% $\sqrt{2}<\ell<2\sqrt{2}$.  Per valori $l\geq 2\sqrt{2}$ il quadrato
% pi\`{u} grande, comunque ruotato, contiene sempre il quadrato di lato
% 1, e pertanto l'area dell'intersezione non dipende nuovamente
% dall'angolo di rotazione.
% 
% Possiamo pertanto concludere che l'area dipende dall'angolo se e solo 
% se $l\geq 2\sqrt{2}$. 
}

\renewcommand{\Commenti}{
}